\section{A degree-twelve logarithm at the old coordinate cost}

These changes start from the delivered 403-coordinate construction, whose
builder has SHA256
\texttt{ca171bf25610cb7da003ed58b4c4d43799a67fd4b5d6ca1b8fbac74a2dcc7e0e}.
The field remains $K=\mathbb Q(\rho)$, $\rho^2=2$, and
$D=(3-\rho)^2$. The inherited global arithmetic inputs remain conditional.

Write
\[
 G_d(w)=\sum_{j=0}^d a_jw^j,\qquad
 a_0=1,\quad a_1=-2/5,\quad
 a_j=\frac{4j-5}{4j+1}a_{j-1}\quad(j\ge2).
\]
The old degree-nine circuit retains the three intermediate field
coordinates $w^2,w^3$, and a middle polynomial before its final logarithm
quotient. A different circuit handles degree twelve with the same number
of coordinates. The pointwise field parameter $w=Df^2$ and its square
$f^2$ are already used as before.

\subsection{Complete a polynomial square}

Choose
\[
 U(w)=w^6+u_5w^5+\cdots+u_1w,\qquad
 R(w)=G_{12}(w)-a_{12}U(w)^2.
\]
Successively match coefficients of degrees $11,10,9,8,7$; each equation
is linear in the next $u_j$, with nonzero coefficient $2a_{12}$.
This determines rational $u_5,\ldots,u_1$ and gives $\deg R\le6$.
The constant coefficient of $U$ is set to zero. Thus this is an exact
polynomial identity over $\mathbb Q$, not an approximation or a choice
of an algebraic square root.

Introduce $w_2=w^2$ and $w_4=w_2^2$. Every polynomial of degree at most
six has a quadratic expression in $1,w,w_2,w_4$, using the basis
\[
 1,\quad w,\quad w_2,\quad ww_2,\quad w_4,
 \quad ww_4,\quad w_2w_4.
\]
Materialize $U$ once by such a quadratic expression. Then
$a_{12}U^2+R$ is also a quadratic expression in the materialized
coordinates. It can be inserted into the old final field quotient.
The final polynomial-output coordinate is eliminated by the same exact
compiler step as before. The three retained intermediate field
coordinates are now $w_2,w_4,U$. The fixed coefficient, if represented
subsequently by a coordinate, does not enter these formulas.

Consequently replacing the four long $G_9$ blocks by $G_{12}$ saves no
core variables and adds none. The isolated full compilation remains
403 coordinates and 371 residuals. The input-only $G_2$ block is retained.
A proposed input upgrade to $G_4$ is not treated as free: the old $G_2$
polynomial output is eliminated, whereas a completed-square $G_4$ needs
another retained field coordinate.

\subsection{Nonvanishing and the consumed local precision}

For $j\ge1$ all $a_j$ are negative, and direct exact addition gives
\[
 G_{12}(1)=\frac{27378487}{130879175}>0.
\]
On every finite real reciprocal chart $0<Df^2<1$. Therefore
$G_{12}(Df^2)\ge G_{12}(1)>0$ at both real embeddings. The old
nonzero logarithm denominators and ratio domains are preserved.
No point chart, sign choice, identity branch, or tangent branch changes
in this logarithm replacement.

The formal-series identity and coefficient bound from the pinned
source, \texttt{manuscripts/quartic\_v10.tex}, lines 9220--9258, give
\[
 \log_E(Q)=z(Q)G_\infty(w),\qquad
 v(w)=4r,\quad r=v(z(Q))\ge1,\qquad
 v(a_j)\ge-\frac{4j+1}{p-1}.
\]
At every unramified consumed prime $p\ge11$, for $j\ge13$ one has
\[
 \frac{4j+1}{p-1}\le\frac{4j+1}{10}
 \le(4j-44)r.
\]
Indeed $10(4j-44)-(4j+1)=36(j-13)+27\ge0$.
Thus every omitted term has valuation at least $44r$, and the bounds
tend to infinity. This proves the relative precision forty-four used
by the selected construction, including the smaller admitted primes.
It does not claim relative precision fifty-two at all those primes.
For example, an omitted coefficient can lose a unit of valuation.

The same formal estimate with first omitted index one gives
$\log_E(Q)/z(Q)\in1+p^{4r}\mathcal O$: $a_1$ is integral for $p\ge11$,
and for $j\ge2$ the inequality
$(4j+1)/10\le4(j-1)r$ suffices. In particular $G_{12}$ has the same
leading-unit and valuation behavior as the exact logarithm.
The existing two-anchor comparison proof therefore uses the new
truncation at order forty-four with no additional point depth or
prime-pattern requirement. The input $G_2$/order-twelve consumer and
the precise finite exceptional set are checked in the separate
round-five arithmetic argument.

\section{A scalar test using only a local approximation to the ratio}

The scalar argument consumes a valuation, so computing the full exact
field ratio is unnecessary. The following replacement saves one
rational coordinate after the full compiler and preserves the actual
scalar divisibility conclusion. It is a transfer of a local predicate,
not a bijection between the old and new norm-tree witness spaces.

Write the fixed point as $P=(x,hj)$ with
\[
 x=h^2,\qquad j^2=x^2-D,\qquad
 h=-1+\frac74\rho.
\]
Retain the actual scalar root $S=(x_S,h_Sj_S)$, its curve equations,
and the shifted point $Q=S-P$. The preceding fused subtraction chart
uses
\[
 B=xx_S+D,\qquad
 Bw=hh_Sj_S-x_Sj,\qquad x_Sf=w^2.
\]
Here $B$ is totally positive, $h_S\ne0$, and $w=h_S\tau$ with
$f=\tau^2=f(Q)$. The deleted intermediate was
\[
 M=-\frac{x+Df}{w},\qquad k_Q=\frac{M-j}{h}.
\]
As before, $k_T=y(T)/x(T)=-1/z(T)$ for a nonidentity finite point.
The shared point is $H=[h_0]P$, and $Q=[a-1]P$.

\subsection{Restore exactly the deleted nonzero restriction}

Deleting the $M$ equation would otherwise allow $w=f=0$, corresponding
to $S=P$. This case must not be admitted. If
$\tau=\tau_0+\tau_1\rho$, then
\[
 f_0=\tau_0^2+2\tau_1^2.
\]
Add the single rational equation $f_0v=1$. It forces $f\ne0$ and
$w\ne0$, and conversely every old admissible point has $f_0>0$ and
the unique extension $v=1/f_0$.

One may now recover $M$ and $k_Q$ by the displayed rational formulas.
The previous complete subtraction proof recovers exactly $Q=S-P$,
including $S=-P$ and excluding $S=P$. No new ordinate branch occurs.
Also $k_Q\ne0$: otherwise its reciprocal square chart would give
$D\tau^4=1$, and hence $7\operatorname N(\tau)^2=1$, impossible
over $\mathbb Q$. This is the same nonvanishing obstruction as before.
The anchor, the two coefficients of $f(Q)$, and all their other
consumers remain unchanged.

\subsection{A uniform two-place trace argument, including the prime 41}

Replace the exact scalar ratio $k_H/k_Q$ by the proxy
\[
 T=k_Hw,
\]
and test only its rational coefficient $T_0$ at the chosen scalar order.
At an odd consumed prime the first four rational norm-tree leaves still
put $H$ and $Q$ in the formal group at every place above $p$.
No conclusion from the proxy is used to obtain this premise.
The inherited logarithm homomorphism and local unit estimate give
\[
 \frac{k_H}{k_Q}=\frac{a-1}{h_0}\,u,
 \qquad u\in1+p^4\mathcal O
\]
at every such place, with the same global field element $u$.

Since $Q$ is formal, the point $S=P+Q$ reduces to the finite integral
point $P$. Thus $h_S$ is integral. Also $\tau\in p\mathcal O$ and
\[
 l_Q=\tau k_Q,\qquad l_Q^2=1-Df^2,
\]
so $l_Q$ is an integral unit and $w\in p\mathcal O$.
The exact recovered chart identity is
\[
 h(wk_Q+h)+Df+jw=0.                         \tag{*}
\]
If $h$ is a unit, this proves $wk_Q\equiv-h\pmod p$.
If $h\in p\mathcal O$, then $x_S\equiv x=h^2\equiv0$, so
$h_S\in p\mathcal O$ and
$wk_Q=h_Sl_Q\equiv0\equiv-h\pmod p$.
This second case is essential at the one place above 41 where $h$ is
not a unit. All fixed coefficients $h,j,x$ are integral outside 2.

Therefore $wk_Q$ is integral and congruent to $-h$ at every place,
including both places in the split case. Multiplication by $u$ preserves
this congruence. Writing
\[
 T=\frac{a-1}{h_0}\,U,\qquad U=u\,wk_Q,
\]
one obtains
\[
 U_0=\frac12\operatorname{Tr}_{K/\mathbb Q}(U)
       \equiv-\frac12\operatorname{Tr}_{K/\mathbb Q}(h)
       =1\pmod p.
\]
Unramifiedness and the two-place integrality assertion justify this
rational trace congruence; $p$ is odd. Consequently
\[
 \boxed{\ v_p(T_0)=v_p((a-1)/h_0).\ }
\]
The scalar norm descent now yields $v_p(a-1)\ge m$ exactly as before.
This works for every positive scalar order $m$, including the selected
order twenty-eight. It introduces no extra exclusion of 41.
Using $-T/h$ instead would have required a separate repair at 41;
the unnormalized proxy deliberately avoids that problem.

\subsection{Completeness and the implemented cost}

At the selected good contact primes, the valuation equality transfers
the old scalar lower bounds directly. At the remaining finite set, the
old preparation permits arbitrary additional depth. The point $S$ tends
to $P$ as $Q$ tends to the identity; $h_S$ remains bounded, while
$w=h_S\tau$ tends to zero. Thus choosing $Q$ sufficiently deeper than
$H$ makes $T_0/r^m$ satisfy any required finite lower bound. Only finitely
many such bounds are altered. The constant left path of the norm tree
is unchanged, so its unit and positivity arguments are preserved.

At a real embedding, $l_Q^2=1-Df^2$ gives $|l_Q|\le1$ and hence
\[
 |T|=|k_H/k_Q|\,|h_Sl_Q|
       \le |h_S|\,|k_H/k_Q|.
\]
The old bounded root chart and real continuity preparation therefore
control the new test after changing finite constants. The inherited
pole-realization and prime-pattern inputs are unchanged. The new
inverse of $f_0$ exists for every old completeness witness.

The actual implementation materializes the single rational quadratic
\[
 T_0=(k_H)_0w_0+2(k_H)_1w_1,
\]
and then its product with $r^{-m}$. This replaces the two rational
coordinates of $M$ and the two of the full field ratio. The extra
inverse of $f_0$ costs one rational coordinate. After accounting for
the old compiler's elimination of the two preliminary scaled-ratio
coordinates, the isolated final saving is one: 402 coordinates and
370 residuals. Every new residual is quadratic.

An initially attractive exact-ratio fusion was rejected by measurement.
It removed $M$, but its cubic numerator forced two preliminary field
components to remain materialized; with the necessary nonzero guard
it compiled to 404, not fewer than 403. The accepted local proxy is
different and has its own soundness and completeness proof above.

\subsection{Finite verification boundary}

The independent source checker reconstructs the point, ratio, and
logarithm equations from the raw IR without importing the producer.
It verifies the completed-square identity, its actual quadratic
coordinate degree, the scalar proxy and inverse guard, and the previous
orientation and exceptional cases. An additional exact split-prime
example uses $Q=[26]P$, $H=[52]P$ at 41: $P$ has reduction orders
13 and 2 at the two places. The proxy's rational unit has residue one,
including the place where $h$ is nonunit. This example supplements,
and does not replace, the uniform trace proof.

The selected integrated source has 401 total coordinates after combining
this one-coordinate scalar saving with the independently proved
one-coordinate scale simplification. The separate arithmetic and detector
proofs justify its order-forty-four tests, finite exceptional set, scalar
order twenty-eight, and anchor exponent thirteen.
